% answers_short.tex —— 1987 年·广东卷·理工农医类 · 简短答案速查
% 由 tools/make_short_answers.py 自 answers.tex 的 \gksolhead/\gkanswer 提取，
% 并与 1987/1987_数学真题_广东理_答案要点.md 逐题比对后生成，勿手改（改源 MD 后重跑）。

\begin{gknotice}
\textbf{使用说明}：本节只给最终答案便于快速核对，完整推导见解析版。答案为编写者独立推导，\textbf{非官方标准答案}。
\end{gknotice}

\gkpart{大题}
\begin{itemize}[leftmargin=2.6em]
  \item[\textbf{1.}] (1) C；(2) D；(3) B；(4) C；(5) A；(6) B；(7) B；(8) D；(9) A；(10) B；(11) A；(12) D；(13) A；(14) B；(15) C；(16) D；(17) 固定源答案表记为 C；但按可见题面条件不能推出 $\alpha\parallel\beta$，见问题台账和解析正文；(18) B；(19) B；(20) D；(21) D；(22) 固定源答案表记为 C；但 A—D 图形选项缺失，字母无法独立核对；(23) D；(24) A；(25) B
  \item[\textbf{2.}] (1) $2$；(2) $3$；(3) 如果 $a,b,c>0$，那么 $\sqrt[3]{abc}\leq\frac{a+b+c}{3}$；(4) $\displaystyle\frac{x^2}{4}+\frac{(y-1)^2}{8}=1$；(5) $-3$
  \item[\textbf{3.}] (1) $BC\perp A_1B,\ BC\perp A_1A,\ A_1A\perp A_1B$（独立推导）；(2) $\sqrt2\text{ cm}$（独立推导）
  \item[\textbf{4.}] (1) $\operatorname{tg}(\alpha+\beta)=\dfrac{24}{7}$；(2) $z^2+u^2+zu=0$
  \item[\textbf{5.}] (1) $4x_1x_2=p^2$；(2) 直线 $l$ 不是弦 $CD$ 的垂直平分线
  \item[\textbf{6.}] (1) $\displaystyle S_n=\frac{a\lg|a|}{(1+a)^2}\left[1+(-1)^{n+1}(1+n+na)a^n\right]$（$a\ne-1$）；(2) 存在。令 $\displaystyle K=\left\lfloor\frac{a^2}{1-a^2}\right\rfloor$，取 $M=2K+2$，则对任意自然数 $n$ 都有 $b_n\leq b_M$
\end{itemize}
